% Options for packages loaded elsewhere
% Options for packages loaded elsewhere
\PassOptionsToPackage{unicode}{hyperref}
\PassOptionsToPackage{hyphens}{url}
\PassOptionsToPackage{dvipsnames,svgnames,x11names}{xcolor}
%
\documentclass[
  10pt,
  letterpaper,
  DIV=11,
  numbers=noendperiod]{scrartcl}
\usepackage{xcolor}
\usepackage[margin=1in]{geometry}
\usepackage{amsmath,amssymb}
\setcounter{secnumdepth}{5}
\usepackage{iftex}
\ifPDFTeX
  \usepackage[T1]{fontenc}
  \usepackage[utf8]{inputenc}
  \usepackage{textcomp} % provide euro and other symbols
\else % if luatex or xetex
  \usepackage{unicode-math} % this also loads fontspec
  \defaultfontfeatures{Scale=MatchLowercase}
  \defaultfontfeatures[\rmfamily]{Ligatures=TeX,Scale=1}
\fi
\usepackage{lmodern}
\ifPDFTeX\else
  % xetex/luatex font selection
\fi
% Use upquote if available, for straight quotes in verbatim environments
\IfFileExists{upquote.sty}{\usepackage{upquote}}{}
\IfFileExists{microtype.sty}{% use microtype if available
  \usepackage[]{microtype}
  \UseMicrotypeSet[protrusion]{basicmath} % disable protrusion for tt fonts
}{}
\makeatletter
\@ifundefined{KOMAClassName}{% if non-KOMA class
  \IfFileExists{parskip.sty}{%
    \usepackage{parskip}
  }{% else
    \setlength{\parindent}{0pt}
    \setlength{\parskip}{6pt plus 2pt minus 1pt}}
}{% if KOMA class
  \KOMAoptions{parskip=half}}
\makeatother
% Make \paragraph and \subparagraph free-standing
\makeatletter
\ifx\paragraph\undefined\else
  \let\oldparagraph\paragraph
  \renewcommand{\paragraph}{
    \@ifstar
      \xxxParagraphStar
      \xxxParagraphNoStar
  }
  \newcommand{\xxxParagraphStar}[1]{\oldparagraph*{#1}\mbox{}}
  \newcommand{\xxxParagraphNoStar}[1]{\oldparagraph{#1}\mbox{}}
\fi
\ifx\subparagraph\undefined\else
  \let\oldsubparagraph\subparagraph
  \renewcommand{\subparagraph}{
    \@ifstar
      \xxxSubParagraphStar
      \xxxSubParagraphNoStar
  }
  \newcommand{\xxxSubParagraphStar}[1]{\oldsubparagraph*{#1}\mbox{}}
  \newcommand{\xxxSubParagraphNoStar}[1]{\oldsubparagraph{#1}\mbox{}}
\fi
\makeatother

\usepackage{color}
\usepackage{fancyvrb}
\newcommand{\VerbBar}{|}
\newcommand{\VERB}{\Verb[commandchars=\\\{\}]}
\DefineVerbatimEnvironment{Highlighting}{Verbatim}{commandchars=\\\{\}}
% Add ',fontsize=\small' for more characters per line
\usepackage{framed}
\definecolor{shadecolor}{RGB}{241,243,245}
\newenvironment{Shaded}{\begin{snugshade}}{\end{snugshade}}
\newcommand{\AlertTok}[1]{\textcolor[rgb]{0.68,0.00,0.00}{#1}}
\newcommand{\AnnotationTok}[1]{\textcolor[rgb]{0.37,0.37,0.37}{#1}}
\newcommand{\AttributeTok}[1]{\textcolor[rgb]{0.40,0.45,0.13}{#1}}
\newcommand{\BaseNTok}[1]{\textcolor[rgb]{0.68,0.00,0.00}{#1}}
\newcommand{\BuiltInTok}[1]{\textcolor[rgb]{0.00,0.23,0.31}{#1}}
\newcommand{\CharTok}[1]{\textcolor[rgb]{0.13,0.47,0.30}{#1}}
\newcommand{\CommentTok}[1]{\textcolor[rgb]{0.37,0.37,0.37}{#1}}
\newcommand{\CommentVarTok}[1]{\textcolor[rgb]{0.37,0.37,0.37}{\textit{#1}}}
\newcommand{\ConstantTok}[1]{\textcolor[rgb]{0.56,0.35,0.01}{#1}}
\newcommand{\ControlFlowTok}[1]{\textcolor[rgb]{0.00,0.23,0.31}{\textbf{#1}}}
\newcommand{\DataTypeTok}[1]{\textcolor[rgb]{0.68,0.00,0.00}{#1}}
\newcommand{\DecValTok}[1]{\textcolor[rgb]{0.68,0.00,0.00}{#1}}
\newcommand{\DocumentationTok}[1]{\textcolor[rgb]{0.37,0.37,0.37}{\textit{#1}}}
\newcommand{\ErrorTok}[1]{\textcolor[rgb]{0.68,0.00,0.00}{#1}}
\newcommand{\ExtensionTok}[1]{\textcolor[rgb]{0.00,0.23,0.31}{#1}}
\newcommand{\FloatTok}[1]{\textcolor[rgb]{0.68,0.00,0.00}{#1}}
\newcommand{\FunctionTok}[1]{\textcolor[rgb]{0.28,0.35,0.67}{#1}}
\newcommand{\ImportTok}[1]{\textcolor[rgb]{0.00,0.46,0.62}{#1}}
\newcommand{\InformationTok}[1]{\textcolor[rgb]{0.37,0.37,0.37}{#1}}
\newcommand{\KeywordTok}[1]{\textcolor[rgb]{0.00,0.23,0.31}{\textbf{#1}}}
\newcommand{\NormalTok}[1]{\textcolor[rgb]{0.00,0.23,0.31}{#1}}
\newcommand{\OperatorTok}[1]{\textcolor[rgb]{0.37,0.37,0.37}{#1}}
\newcommand{\OtherTok}[1]{\textcolor[rgb]{0.00,0.23,0.31}{#1}}
\newcommand{\PreprocessorTok}[1]{\textcolor[rgb]{0.68,0.00,0.00}{#1}}
\newcommand{\RegionMarkerTok}[1]{\textcolor[rgb]{0.00,0.23,0.31}{#1}}
\newcommand{\SpecialCharTok}[1]{\textcolor[rgb]{0.37,0.37,0.37}{#1}}
\newcommand{\SpecialStringTok}[1]{\textcolor[rgb]{0.13,0.47,0.30}{#1}}
\newcommand{\StringTok}[1]{\textcolor[rgb]{0.13,0.47,0.30}{#1}}
\newcommand{\VariableTok}[1]{\textcolor[rgb]{0.07,0.07,0.07}{#1}}
\newcommand{\VerbatimStringTok}[1]{\textcolor[rgb]{0.13,0.47,0.30}{#1}}
\newcommand{\WarningTok}[1]{\textcolor[rgb]{0.37,0.37,0.37}{\textit{#1}}}

\usepackage{longtable,booktabs,array}
\usepackage{calc} % for calculating minipage widths
% Correct order of tables after \paragraph or \subparagraph
\usepackage{etoolbox}
\makeatletter
\patchcmd\longtable{\par}{\if@noskipsec\mbox{}\fi\par}{}{}
\makeatother
% Allow footnotes in longtable head/foot
\IfFileExists{footnotehyper.sty}{\usepackage{footnotehyper}}{\usepackage{footnote}}
\makesavenoteenv{longtable}
\usepackage{graphicx}
\makeatletter
\newsavebox\pandoc@box
\newcommand*\pandocbounded[1]{% scales image to fit in text height/width
  \sbox\pandoc@box{#1}%
  \Gscale@div\@tempa{\textheight}{\dimexpr\ht\pandoc@box+\dp\pandoc@box\relax}%
  \Gscale@div\@tempb{\linewidth}{\wd\pandoc@box}%
  \ifdim\@tempb\p@<\@tempa\p@\let\@tempa\@tempb\fi% select the smaller of both
  \ifdim\@tempa\p@<\p@\scalebox{\@tempa}{\usebox\pandoc@box}%
  \else\usebox{\pandoc@box}%
  \fi%
}
% Set default figure placement to htbp
\def\fps@figure{htbp}
\makeatother





\setlength{\emergencystretch}{3em} % prevent overfull lines

\providecommand{\tightlist}{%
  \setlength{\itemsep}{0pt}\setlength{\parskip}{0pt}}



 


\KOMAoption{captions}{tableheading}
\makeatletter
\@ifpackageloaded{caption}{}{\usepackage{caption}}
\AtBeginDocument{%
\ifdefined\contentsname
  \renewcommand*\contentsname{Table of contents}
\else
  \newcommand\contentsname{Table of contents}
\fi
\ifdefined\listfigurename
  \renewcommand*\listfigurename{List of Figures}
\else
  \newcommand\listfigurename{List of Figures}
\fi
\ifdefined\listtablename
  \renewcommand*\listtablename{List of Tables}
\else
  \newcommand\listtablename{List of Tables}
\fi
\ifdefined\figurename
  \renewcommand*\figurename{Figure}
\else
  \newcommand\figurename{Figure}
\fi
\ifdefined\tablename
  \renewcommand*\tablename{Table}
\else
  \newcommand\tablename{Table}
\fi
}
\@ifpackageloaded{float}{}{\usepackage{float}}
\floatstyle{ruled}
\@ifundefined{c@chapter}{\newfloat{codelisting}{h}{lop}}{\newfloat{codelisting}{h}{lop}[chapter]}
\floatname{codelisting}{Listing}
\newcommand*\listoflistings{\listof{codelisting}{List of Listings}}
\makeatother
\makeatletter
\makeatother
\makeatletter
\@ifpackageloaded{caption}{}{\usepackage{caption}}
\@ifpackageloaded{subcaption}{}{\usepackage{subcaption}}
\makeatother
\usepackage{bookmark}
\IfFileExists{xurl.sty}{\usepackage{xurl}}{} % add URL line breaks if available
\urlstyle{same}
\hypersetup{
  pdftitle={Contraction Theory Meets Tangent Spaces: Stability, Optimality, and Scope Limits},
  pdfauthor={Dieter Olson},
  colorlinks=true,
  linkcolor={blue},
  filecolor={Maroon},
  citecolor={Blue},
  urlcolor={Blue},
  pdfcreator={LaTeX via pandoc}}


\title{Contraction Theory Meets Tangent Spaces: Stability, Optimality,
and Scope Limits}
\usepackage{etoolbox}
\makeatletter
\providecommand{\subtitle}[1]{% add subtitle to \maketitle
  \apptocmd{\@title}{\par {\large #1 \par}}{}{}
}
\makeatother
\subtitle{What a Differential Stability Certificate Establishes}
\author{Dieter Olson}
\date{2026-01-18}
\begin{document}
\maketitle

\renewcommand*\contentsname{Contents}
{
\hypersetup{linkcolor=}
\setcounter{tocdepth}{1}
\tableofcontents
}

\clearpage

\section*{Abstract}\label{abstract}
\addcontentsline{toc}{section}{Abstract}

Contraction analysis asks whether neighboring solutions approach one
another under a specified dynamics model and input policy. Trajectory
optimization asks which admissible motion minimizes a declared
objective. Their common use of differential dynamics creates a useful
connection: under explicit definiteness and boundedness conditions, a
linear-quadratic value matrix also certifies contraction of the
associated closed loop. Neither optimality nor a positive matrix alone
supplies that certificate.

This article derives the connection, distinguishes continuous and
discrete time, and develops counterexamples to stronger interpretations.
It then treats metric geometry, numerical verification, musculoskeletal
actuation, stochastic variability, task-space control, and contact
events. The golf application connects geometry and stored energy to
feasible corrections and the delivered impact state. Its empirical
questions remain separate from the mathematical identities.

\section{Reading the Physical
Argument}\label{reading-the-physical-argument}

A repeatable swing need not erase every state difference. It may
tolerate variation in directions that have little influence on club
delivery while controlling a small number of consequential errors. A
contraction metric is one way to test a stronger model claim: all
differential errors measured by that metric decrease throughout a
specified region and time interval.

The metric changes how error is measured; feedback and mechanics change
how error evolves. Calling a metric curved does not make it a physical
restoring force. Earlier work, geometry, inertia, grip compliance,
activation history, and ground contact determine the state and the
available responses. A stability calculation must include those
mechanisms in its dynamics before drawing conclusions about their
effects.

\section{Differential Dynamics and the Correct Metric
Test}\label{differential-dynamics-and-the-correct-metric-test}

Let the smooth closed-loop dynamics be
\begin{equation}\phantomsection\label{eq-ctu-dynamics}{
\dot x=F(x,t)=f(x,\pi(x,t),t),\qquad A_{\mathrm{cl}}=D_xF.
}\end{equation} Here \(\pi\) is the actual input policy. Differentiating
while holding time fixed gives the exact variational equation
\begin{equation}\phantomsection\label{eq-virtual-displacement}{
\delta\dot x=A_{\mathrm{cl}}(x,t)\delta x.
}\end{equation} For an open-loop plant with independently perturbed
input, the equation is instead \(\delta\dot x=f_x\delta x+f_u\delta u\).
Closing the loop adds \(\delta u=\pi_x\delta x\), so
\(A_{\mathrm{cl}}=f_x+f_u\pi_x\). Differentiating the plant with input
held fixed does not test feedback stability.

Choose a smooth symmetric metric \(M(x,t)\) and define differential
squared length \(V=\delta x^\top M\delta x\). Its derivative is
\begin{equation}\phantomsection\label{eq-distance-derivative}{
\begin{aligned}
\dot V&=\delta x^\top\mathcal H\delta x,\\
\mathcal H&=A_{\mathrm{cl}}^\top M+MA_{\mathrm{cl}}+\dot M,\\
\dot M&=\partial_tM+\sum_iF_i\partial_{x_i}M.
\end{aligned}
}\end{equation} The material derivative uses the same closed-loop flow.
A sequence of matrices evaluated along one trajectory does not define
the spatial derivatives of a metric in its neighborhood.

The contraction inequality with length-decay rate \(\lambda>0\) is
\begin{equation}\phantomsection\label{eq-contraction-condition}{
\mathcal H\preceq-2\lambda M.
}\end{equation} Its normalized symmetric matrix is
\begin{equation}\phantomsection\label{eq-generalized-jacobian}{
M^{-1/2}\mathcal H M^{-1/2}\preceq-2\lambda I.
}\end{equation} The multiplication order matters. With \[
A=\begin{bmatrix}-1&3\\0&-2\end{bmatrix},\qquad
M=\begin{bmatrix}1/2&1/2\\1/2&1\end{bmatrix},
\] the correct \(A^\top M+MA=-I\). The different expression
\(AM+MA^\top\) has a positive eigenvalue approximately \(2.354\) and
tests a different object.

For a constant LTI matrix and constant metric this specializes to
\begin{equation}\phantomsection\label{eq-linear-contraction}{
A^\top M+MA\preceq-2\lambda M,\qquad M\succ0.
}\end{equation}

If \(M=\Theta^\top\Theta\) with invertible \(\Theta\), differential
components \(\delta z=\Theta\delta x\) obey
\begin{equation}\phantomsection\label{eq-ctu-generalized-jacobian-def}{
\delta\dot z=\mathcal F\delta z,\qquad
\mathcal F=(\dot\Theta+\Theta A_{\mathrm{cl}})\Theta^{-1}.
}\end{equation} Then \(\operatorname{sym}\mathcal F\preceq-\lambda I\)
is equivalent to the metric inequality. The differential transformation
need not integrate to a global coordinate map \(z(x)\).

\subsection{From Differential Length to Finite
Separation}\label{from-differential-length-to-finite-separation}

Integrating the differential inequality gives
\begin{equation}\phantomsection\label{eq-exponential-convergence}{
\|\delta x(t)\|_{M(x(t),t)}
\leq e^{-\lambda(t-t_0)}\|\delta x(t_0)\|_{M(x(t_0),t_0)}.
}\end{equation} For finite separation, connect the initial states by a
smooth path and evolve every point of that path under the same flow. If
this path of solutions remains inside the certified region, its metric
length shrinks. Taking the infimum over admissible connecting paths
yields a corresponding bound on Riemannian distance. A finite difference
\(x_2-x_1\) cannot simply be substituted for \(\delta x\) in the exact
differential equation of a nonlinear system.

State existence, path containment, and metric bounds belong to the
conclusion. Uniform bounds \(0<m_-I\preceq M\preceq m_+I\) permit
comparison with Euclidean errors, including the conditioning factor
\(\sqrt{m_+/m_-}\). An invariant region or an invariant tube about a
reference can provide the required containment. A condition checked only
at the reference does not establish the tube. These distinctions follow
the path-based formulation in
\href{https://web.mit.edu/nsl/www/preprints/contraction.pdf}{Lohmiller
and Slotine}.

Contraction shrinks every differential length in the chosen metric.
Volume shrinkage is weaker: \(\dot x=\operatorname{diag}(1,-2)x\)
shrinks area while one direction expands. Convergence between solutions
also does not imply bounded motion: \(\dot x=-x+e^{2t}\) has contracting
differences while its solutions grow without bound. For an autonomous
system, strict full-state contraction cannot support a nonconstant
periodic orbit in the same certified domain, because phase-shifted
solutions would have to converge at equal time. Orbital or transverse
contraction is a different statement.

\subsection{Disturbances and
Interconnections}\label{disturbances-and-interconnections}

Suppose one trajectory receives an added disturbance whose norm in the
certified metric is bounded by \(\bar d\). Under the same region and
regularity conditions, a comparison inequality gives \[
d_M(t)\leq e^{-\lambda(t-t_0)}d_M(t_0)
+\frac{\bar d}{\lambda}\left(1-e^{-\lambda(t-t_0)}\right).
\] Persistent disturbance generally produces a nonzero error bound. It
does not guarantee convergence to the undisturbed trajectory, and the
perturbed trajectory must remain in the region where the bound applies.

Independent contracting subsystems driven by the same prescribed signal
admit a block-diagonal contraction metric. Mechanically coupled joints
are not independent subsystems. Feedback interconnection adds
off-diagonal Jacobian blocks whose symmetric contribution must be
bounded. A bounded hierarchical coupling can sometimes be handled by
rescaling a composite metric; arbitrary mutual coupling cannot be
ignored. For example, two isolated equations \(\dot x_i=-x_i\) are
stable, but coupling them as \(\dot x_1=-x_1+2x_2\),
\(\dot x_2=2x_1-x_2\) creates an expanding eigenvalue \(1\).

\section{Riccati Value Functions and
Contraction}\label{riccati-value-functions-and-contraction}

\subsection{The Continuous-Time
Identity}\label{the-continuous-time-identity}

Consider a linear time-varying perturbation system
\(\dot\xi=A(t)\xi+B(t)v\) and the quadratic cost \[
J=\xi(T)^\top S_T\xi(T)
+\int_{t_0}^T\left(\xi^\top Q\xi+2\xi^\top Nv+v^\top Rv\right)dt,
\qquad R\succ0.
\] Assume the cost block is positive semidefinite and the Riccati
solution exists. Completing the square in dynamic programming gives
\begin{equation}\phantomsection\label{eq-differential-riccati}{
-\dot S=A^\top S+SA-(SB+N)R^{-1}(B^\top S+N^\top)+Q,
\qquad
K=R^{-1}(B^\top S+N^\top).
}\end{equation} For the minimizing policy \(v=-K\xi\), let \(A_c=A-BK\)
and \[
W=Q-NK-K^\top N^\top+K^\top RK.
\] Substitution yields the exact identity
\begin{equation}\phantomsection\label{eq-closed-loop-lyapunov}{
\dot S+A_c^\top S+SA_c=-W.
}\end{equation} Thus \(S\) is a contraction metric on an interval if it
is uniformly positive definite and bounded there and
\(W\succeq2\lambda S\) holds uniformly. For zero cross cost,
\(W=Q+K^\top RK\). The open-loop ARE instead gives
\(A^\top S+SA=SBR^{-1}B^\top S-Q\); the negative feedback contribution
appears only after forming \(A_c\).

For infinite-horizon LTI LQR, standard stabilizability and detectability
assumptions with \(Q\succeq0\), \(R\succ0\) yield a stabilizing
positive-semidefinite solution. Positive definiteness and a strict decay
bound in that particular matrix require their own checks. For example,
\(A=-1\), \(B=1\), \(Q=0\), \(R=1\) has optimal cost \(S=0\) and a
stable uncontrolled plant, but zero is not a Riemannian metric. With
\(Q\succ0\) the usual stabilizing solution is positive definite. The
\href{https://underactuated.mit.edu/lqr.html}{MIT LQR treatment}
develops the finite- and infinite-horizon optimization problems.

\subsection{Finite-Horizon Decrease Is Not an Asymptotic
Certificate}\label{finite-horizon-decrease-is-not-an-asymptotic-certificate}

A terminal weight may vanish. In the scalar uncontrolled system
\(\dot\xi=\xi\) with running cost \(\xi^2\) and zero terminal cost, \[
S(t)=\frac{e^{2(T-t)}-1}{2},\qquad \dot S+2S=-1.
\] The value \(S(t)\xi(t)^2\) decreases as the remaining horizon
disappears, while \(|\xi(t)|\) grows. At \(T\), \(S(T)=0\). There is no
uniformly positive metric through the endpoint and no claim about
behavior after the horizon. This example isolates why value decrease and
physical error decay must be distinguished.

In a nonlinear tracking problem, linearization along a nominal
trajectory supplies a variational candidate. To extend its result to a
finite neighborhood, verify the closed-loop derivative and metric
inequality throughout a tube, with strict margin, model residual bounds,
feasible controls, and containment. Recomputing LQR at sample points
does not supply these conditions automatically.

\subsection{Discrete-Time Dynamics Require a Different
Test}\label{discrete-time-dynamics-require-a-different-test}

For a discrete closed loop \(x_{k+1}=F_k(x_k)\) with differential map
\(A_{c,k}=D F_k\), the contraction condition is
\begin{equation}\phantomsection\label{eq-discrete-contraction}{
A_{c,k}^\top M_{k+1}(F_k(x))A_{c,k}\preceq\rho M_k(x),\qquad0<\rho<1.
}\end{equation} The metric is evaluated at the next state on the left.
At sample interval \(\Delta t\), the length-decay rate is
\(-\log(\rho)/(2\Delta t)\). A continuous Lyapunov inequality applied to
the Jacobian of a discrete integrator is not this condition.

For zero-cross-cost discrete LQR, the Riccati identity is
\begin{equation}\phantomsection\label{eq-riccati-substitution}{
(A_k-B_kK_k)^\top S_{k+1}(A_k-B_kK_k)
=S_k-Q_k-K_k^\top R_kK_k.
}\end{equation} If \(S_k\) has uniform positive bounds and
\(Q_k+K_k^\top R_kK_k\succeq(1-\rho)S_k\), it supplies the desired
contraction inequality. The normalized matrix inequality, not the
spectral radius of each update, establishes the one-step rate.

For example, \[
F_1=\begin{bmatrix}0.5&2\\0&0.5\end{bmatrix},\qquad F_2=F_1^\top.
\] Both matrices have eigenvalues \(0.5\), yet \(F_2F_1\) has an
eigenvalue approximately \(4.486\). Alternation amplifies errors.
Time-varying stability depends on products or a valid changing metric,
not independent eigenvalue tests. Even in an LTI system, nonnormal
transient amplification separates an asymptotic spectral rate from
monotone decay in a chosen norm.

\subsection{Worked Double-Integrator
Certificate}\label{worked-double-integrator-certificate}

Take \[
A=\begin{bmatrix}0&1\\0&0\end{bmatrix},\quad
B=\begin{bmatrix}0\\1\end{bmatrix},\quad
Q=\operatorname{diag}(16,1),\quad R=1.
\] The stabilizing ARE solution and feedback are
\begin{equation}\phantomsection\label{eq-double-integrator-riccati}{
S=\begin{bmatrix}12&4\\4&3\end{bmatrix},\qquad K=\begin{bmatrix}4&3\end{bmatrix}.
}\end{equation} They give eigenvalues \(-1.5\pm(\sqrt7/2)i\). In the
specific metric \(S\), however, the largest certified uniform
length-decay rate is \[
\lambda_S=\frac12\lambda_{\min}\left(S^{-1/2}(Q+K^\top K)S^{-1/2}\right)
\approx1.276393\ \mathrm{s}^{-1}.
\] The other normalized half-eigenvalue is approximately \(1.723607\).
Neither the eigenvalue real part \(1.5\) nor a requested design value
\(2\) can replace this calculation. Rates of exponential decay have
units inverse seconds; they are not oscillation frequencies measured in
cycles per second.

LQR minimizes its declared cost, not contraction rate over all feedback
policies. Changing the cost weights changes that compromise; a scalar
weight ratio cannot prescribe arbitrary multivariable poles. To design a
rate directly for an unconstrained LTI plant, set \(X=M^{-1}\) and
\(Y=KX\). Congruence transforms the state-feedback inequality into \[
AX+XA^\top-BY-Y^\top B^\top+2\lambda X\preceq0,
\qquad X\succ0.
\] For fixed \(\lambda\), this is linear in \(X,Y\). Add normalization
and any required actuator or region constraints, and recover
\(K=YX^{-1}\). This feasibility problem is not automatically an
inverse-LQR solution, and imposing input bounds over a state region
requires additional constraints.

\section{What the Geometry Does and Does Not
Say}\label{what-the-geometry-does-and-does-not-say}

\subsection{Metric Tensors and
Coordinates}\label{metric-tensors-and-coordinates}

For a diffeomorphism \(z=\phi(x)\), write \(T=D\phi\). Since
\(\delta z=T\delta x\), equal differential lengths require \[
M_z=T^{-\top}M_xT^{-1}.
\] The expression \(J^\top M_yJ\) is instead the pullback of a metric
from a target through a map \(y=h(x)\) with \(J=Dh\). These are
compatible tensor rules in opposite directions. Coordinate invariance of
the inequality does not make arbitrary choices of metric equivalent or
make their Euclidean conditioning identical.

The bilinear metric and task pullback can be written explicitly as
\begin{equation}\phantomsection\label{eq-metric-tensor}{
g_x(\xi,\zeta)=\xi^\top M_x\zeta,
}\end{equation}
\begin{equation}\phantomsection\label{eq-pullback-metric}{
h^*M_y=J^\top M_yJ.
}\end{equation}

A Riccati matrix depending on time alone has no spatial derivatives at
each fixed time. Its spatial Riemann curvature is zero in these
coordinates. LQR trajectories are solutions of controlled dynamics and
an optimization problem; they are not generally geodesics of that
matrix. For instance, \(\dot x=-x\) contracts in the flat Euclidean
metric, without curvature or a funnel-shaped physical surface.

For a metric \(g\), the Levi--Civita coefficients are
\begin{equation}\phantomsection\label{eq-christoffel}{
\Gamma^i_{jk}=\tfrac12g^{i\ell}(\partial_jg_{\ell k}+\partial_kg_{\ell j}-\partial_\ell g_{jk}).
}\end{equation} They vanish for a spatially constant metric in these
coordinates. This fact is compatible with a nonzero contraction rate of
a dissipative controlled flow.

The variational equation in coordinates uses ordinary time derivatives.
With a chosen torsion-free connection it can be written
\(\nabla_t\delta x=\nabla_{\delta x}F\); replacing only the left
derivative with a covariant derivative while keeping an ordinary
Jacobian on the right omits connection terms. This geometric rewriting
does not turn the dynamics into an Euler--Lagrange optimality condition.
Differentiating an ODE and varying a dynamics-constrained objective are
separate operations.

In components, the covariant identity is
\begin{equation}\phantomsection\label{eq-tangent-bundle-dynamics}{
\frac{D\delta x^i}{dt}=\left(\partial_jF^i+\Gamma^i_{kj}F^k\right)\delta x^j.
}\end{equation}

\subsection{Task Metrics, Inertia, and
Stiffness}\label{task-metrics-inertia-and-stiffness}

For a task \(y=h(q)\), \(J^\top W_yJ\) measures local task displacement.
It is positive definite only if \(J\) has full column rank when
\(W_y\succ0\). Redundancy or singularity can leave task-invisible
directions. With a straight two-link unit-length arm, \[
J=\begin{bmatrix}0&0\\2&1\end{bmatrix},\qquad
J^\top J=\begin{bmatrix}4&2\\2&1\end{bmatrix},
\] which is singular. It equals the generalized kinetic metric of a unit
point mass at the endpoint with massless links. It is not the mass
matrix of two unit-mass links. For two uniform unit rods the actual
inertia is
\(\left[\begin{smallmatrix}8/3&5/6\\5/6&1/3\end{smallmatrix}\right]\),
which is positive definite at this configuration.

Mechanical inertia acts on configuration velocities. A contraction
metric for state \((q,\dot q)\) acts on both position and velocity
perturbations and generally includes cross blocks. Stiffness acts
between displacement and force; it is not automatically a full-state
contraction metric. Likewise, \(BR^{-1}B^\top\) is a contravariant input
map that is often singular in underactuated mechanics, not automatically
a positive-definite covariant metric. Fisher information and value
Hessians can also be degenerate. Every proposed metric needs a domain,
units or nondimensionalization, rank check, and differential inequality.

Minimizing \(\operatorname{tr}M\) without a fixed lower bound is
degenerate because the homogeneous contraction inequality permits
\(\epsilon M\) for arbitrarily small \(\epsilon>0\). Scaling the ruler
does not enlarge a physical basin. With \(I\preceq M\preceq\kappa I\), a
fixed linear closed loop and fixed rate produce a meaningful
semidefinite feasibility problem. A nonlinear metric search must specify
a function class, derivative constraints, a domain, and how a continuum
inequality is verified. Solver complexity depends on those choices; no
universal \(O(n^4)\) runtime follows from naming the problem an SDP.

For the fixed constant closed loop, the normalized feasibility
constraints are
\begin{equation}\phantomsection\label{eq-contraction-lmi}{
I\preceq M\preceq\kappa I,\qquad A_c^\top M+MA_c+2\lambda M\preceq0.
}\end{equation}

\section{Combining Trajectory Optimization With Stability
Verification}\label{combining-trajectory-optimization-with-stability-verification}

\subsection{What DDP Actually Expands}\label{what-ddp-actually-expands}

For a discrete dynamics map \(x^+=f_d(x,u)\) and running cost \(\ell\),
define the local action value \(\mathcal Q=\ell+V^+(f_d(x,u))\). Along a
feasible nominal trajectory, write \(A=f_{d,x}\), \(B=f_{d,u}\),
\(p=V_x^+\), and \(P=V_{xx}^+\). The second derivatives include \[
\begin{aligned}
\mathcal Q_{xx}&=\ell_{xx}+A^\top PA+\sum_i p_i f_{d,xx}^{,i},\\
\mathcal Q_{ux}&=\ell_{ux}+B^\top PA+\sum_i p_i f_{d,ux}^{,i},\\
\mathcal Q_{uu}&=\ell_{uu}+B^\top PB+\sum_i p_i f_{d,uu}^{,i}.
\end{aligned}
\] DDP retains these dynamics-Hessian contractions. The iLQR
approximation omits them while retaining the relevant cost Hessians and
cross terms. In either case, a general optimization step includes a
feedforward correction as well as feedback: \[
\delta u=k+L\delta x,\qquad
k=-\mathcal Q_{uu}^{-1}\mathcal Q_u,\quad
L=-\mathcal Q_{uu}^{-1}\mathcal Q_{ux},
\] when the local control Hessian is positive definite and the step is
unconstrained. Regularization, input bounds, line search, and defects in
an infeasible nominal rollout change the subproblem. The purely
homogeneous LQR recursion describes a special case, not a complete
nonlinear optimizer. See
\href{https://underactuated.mit.edu/trajopt.html}{MIT trajectory
optimization}.

A value Hessian can be indefinite away from a suitable local minimum.
Even a positive Hessian does not certify the resulting nonlinear
feedback policy. Optimization convergence, trajectory feasibility,
stability, and empirical model validity are distinct checks.

\subsection{A Defensible Contraction
Penalty}\label{a-defensible-contraction-penalty}

For a declared metric and continuous closed-loop policy, define the
normalized defect \[
d(x,t)=\lambda_{\max}\left(M^{-1/2}\mathcal H M^{-1/2}\right)+2\lambda,
\qquad c(x,t)=\max(0,d(x,t))^2.
\] Zero penalty means the tested point satisfies the inequality at the
requested rate. A finite penalty weight in an objective does not force
zero violation. Sampled zero violations also do not certify an unsampled
region. The largest-eigenvalue function is nonsmooth at repeated active
eigenvalues, so an implementation must specify how derivatives or
nonsmooth steps are handled.

If a scalar penalty is added to a cost, its first derivatives modify the
cost gradients and its second derivatives modify the cost Hessians,
including mixed terms. Adding a gradient vector to a Hessian matrix is
dimensionally invalid. If the policy or metric is optimized too, account
for their dependence in those derivatives. Use the continuous condition
for continuous Jacobians and the next-state metric inequality for
discrete maps.

For a candidate policy and metric, a useful verification sequence is:

\begin{enumerate}
\def\labelenumi{\arabic{enumi}.}
\tightlist
\item
  Check the nominal trajectory against the dynamics, contact model, and
  actuator limits.
\item
  Define the metric throughout a proposed tube, with explicit positive
  lower and upper bounds.
\item
  Differentiate the implemented closed loop, including saturation or its
  smooth region, activation states, and metric material derivatives.
\item
  Verify the contraction inequality with numerical margin throughout
  that tube, and establish trajectory containment. A certified Lipschitz
  bound can turn a sampled margin into a neighborhood bound; an
  arbitrary sample spacing cannot.
\item
  Bound model errors and disturbances, inspect contact-event
  sensitivities, and report the resulting finite-error or output bound.
\end{enumerate}

For example, if the largest eigenvalue of the normalized defect matrix
is at most \(-\eta\) at a point and a proven Lipschitz constant is
\(L\), a ball of radius \(r<\eta/L\) retains negativity, provided the
metric and policy remain regular there. This verifies a local
inequality; an invariant tube and time-uniform bounds still need proof.
A finite-horizon capture set under one policy is not automatically an
asymptotic basin of attraction or a control-reachable set under all
policies.

Control contraction metrics provide a separate constructive framework
for differential feedback design. The original
\href{https://arxiv.org/abs/1503.03144}{Manchester--Slotine paper} and
its \href{https://arxiv.org/abs/1711.08128}{integrability amendment}
specify technical conditions for turning differential controls into
finite controls. A pointwise Riccati matrix or a soft penalty should not
be presented as an implementation of those theorems without checking
their hypotheses.

\section{Musculoskeletal Mechanics and
Variability}\label{musculoskeletal-mechanics-and-variability}

\subsection{Force, Activation, and Full Mechanical
State}\label{force-activation-and-full-mechanical-state}

Moment arms map muscle forces to joint torque: \[
\tau=R(q)F_m(q,\dot q,a,z)+\tau_{\mathrm{other}}.
\] Here \(a\) denotes activation and \(z\) any additional tendon,
fatigue, or history state. Multiplying a moment-arm matrix, measured in
length, by dimensionless activation alone does not produce torque. A
simplified constitutive law can be affine in activation at fixed muscle
state; its force-length, force-velocity, tendon, and excitation dynamics
must still be specified. The
\href{https://opensimconfluence.atlassian.net/wiki/spaces/OpenSim/pages/53090590}{OpenSim
activation model} illustrates why excitation and activation occupy
different roles.

A low-dimensional factorization \(a\approx Wc\) is a candidate
description of recorded activation patterns. Its columns live in
muscle-activation space, not directly in the tangent space of the
mechanical state. The resulting input directions involve the muscle
model, moment arms, inertia, and constraints. A low-rank factorization
does not establish that the nervous system selects \(W\) to maximize
contraction, or even that the factors identify unique neural modules.

For mechanics \[
D(q)\ddot q+h(q,v)=B(q)\tau,\qquad v=\dot q,
\] define \(a_q=D^{-1}(B\tau-h)\) as the acceleration function. Its
configuration derivative in direction \(\delta q\) is \[
D_qa_q[\delta q]=D^{-1}\left(D_qB[\delta q]\tau-D_qh[\delta q]-D_qD[\delta q]a_q\right),
\] and \(D_va_q=-D^{-1}D_vh\), \(D_\tau a_q=D^{-1}B\) under the stated
\(q\)-only inertia and input map. The first-order state Jacobian has top
blocks \((0,I)\) and bottom blocks \((D_qa_q,D_va_q)\). Omitting the
derivative of inertia or the velocity dependence of the Coriolis load
changes the predicted perturbation dynamics. Closed loops add reaction
derivatives; muscle-driven models add their internal-state blocks.

\subsection{Impedance Is Part of the
Dynamics}\label{impedance-is-part-of-the-dynamics}

A local endpoint law \(\delta F=-K_y\delta y-C_y\delta\dot y\)
contributes stiffness and damping. The simple pullback \(J^\top K_yJ\)
is the unloaded linearization at a matching reference; differentiating
\(J^\top F\) at nonzero preload also produces a geometric term involving
\(D_qJ\). Along motion, velocity and Jacobian-rate terms matter as well.

Finite stiffness can redirect an error without eliminating a degree of
freedom. A spring without damping stores and returns energy; it does not
make trajectories converge. A contraction test therefore belongs to the
full position/velocity/internal-state dynamics. Co-contraction may alter
impedance and later force production without changing instantaneous net
torque, and it can have metabolic cost. Neither an inferred torque nor a
metric identifies that cost uniquely.

\subsection{Stochastic Errors Need a Noise
Model}\label{stochastic-errors-need-a-noise-model}

For fully observed additive-noise LQ control, \[
dx=(Ax+Bu)dt+G\,dW_t,\qquad \mathbb E[dW_t,dW_t^\top]=I\,dt,
\] the expected quadratic-cost problem has the deterministic Riccati
equation for its state-quadratic part. Noise contributes an additive
value term involving \(\operatorname{tr}(SGG^\top)\). Under \(u=-Kx\),
the covariance evolves with the closed-loop matrix: \[
\dot\Sigma_x=A_c\Sigma_x+\Sigma_xA_c^\top+GG^\top.
\] For \(dx=-x,dt+\sigma,dW_t\) from a known initial state, the variance
is \(\sigma^2(1-e^{-2t})/2\). Mean error decays, but variance approaches
a nonzero floor. Two trajectories driven by exactly the same additive
noise have deterministic difference dynamics; independent noise does not
cancel. These are different notions of stochastic contraction and
synchronization.

Risk-sensitive exponential costs can change the Riccati equation, but
their noise-dependent term follows from the specified criterion,
normalization, and admissibility conditions. It should not be inserted
into ordinary expected-quadratic LQR by analogy. Additive-noise
covariance alone also does not establish that a human must always
increase effort to maintain a particular outcome; the objective,
feedback, signal-dependent noise, and feasible muscle responses need a
model and data.

\section{Task-Space Tracking and
Contact}\label{task-space-tracking-and-contact}

\subsection{Operational-Space Bias Is a Force
Vector}\label{operational-space-bias-is-a-force-vector}

For \(D\ddot q+h=\tau\), task \(y=h_y(q)\), and full-row-rank task
Jacobian \(J\), suppose the input is implemented as
\(\tau=J^\top F+\tau_0\) with \(JD^{-1}\tau_0=0\). Then \[
\Lambda=(JD^{-1}J^\top)^{-1},\qquad
\Lambda\ddot y+\eta=F,\qquad
\eta=\Lambda(JD^{-1}h-\dot Jv).
\] The bias \(\eta\) is a force vector, not a matrix that should be
multiplied by task velocity again. Its state dependence generally cannot
be expressed using only \(y,\dot y\) in a redundant arm. Rank loss,
constraint contacts, torque limits, and imperfect cancellation require
separate treatment.

With exact model cancellation, choose
\begin{equation}\phantomsection\label{eq-task-force}{
F=\eta+\Lambda\left(\ddot y_d-2a\dot e-a^2e\right),\qquad e=y-y_d,
}\end{equation} to obtain \(\ddot e+2a\dot e+a^2e=0\). This is
critically damped. Its solutions contain \((c_1+c_2t)e^{-at}\), so a
uniform exponential bound with exactly rate \(a\) and fixed prefactor
does not hold for arbitrary initial error. Any strictly smaller rate can
be bounded with a suitable prefactor. The full error-state metric has
dimension twice the task dimension; \(\Lambda\) alone cannot serve that
role.

For \(a=1\) and error state \(z=(e,\dot e)\), \[
A_e=\begin{bmatrix}0&1\\-1&-2\end{bmatrix},\qquad
P=\begin{bmatrix}3/2&1/2\\1/2&1/2\end{bmatrix}
\] satisfies \(A_e^\top P+PA_e=-I\). It gives the valid rate
\(1/(2\lambda_{\max}(P))\approx0.292893\ \mathrm{s}^{-1}\) in that
metric. In contrast, the simple energy \((e^2+\dot e^2)/2\) has
derivative \(-2\dot e^2\), which is zero at \(e\ne0,\dot e=0\) and
cannot satisfy a strict instantaneous exponential energy inequality. A
conservative certificate can be correct even when it is below the
asymptotic spectral rate.

\subsection{Saltation and Mode
Changes}\label{saltation-and-mode-changes}

A contact model needs a guard, pre- and post-contact vector fields, and
a reset map. For a smooth non-grazing event \(g(x,t)=0\) with reset
\(x^+=R(x^-,t)\), the first-order sensitivity across the event is \[
\Xi=D_xR+
\frac{F^+-D_xR\,F^- -\partial_tR}
{D_xg\,F^-+\partial_tg}\,D_xg.
\] This is the saltation matrix. It includes event-time shifts as well
as reset derivatives; see the
\href{https://arxiv.org/abs/2306.06862}{saltation review}. The
denominator must be nonzero, and a consistent coordinate representation
and event sequence are required. Grazing, simultaneous contacts, and
changes in effective state dimension need additional analysis.

Across an event, a nonexpansion condition is \[
\Xi^\top M^+\Xi\preceq M^-.
\] For constant linear mode dynamics, a common flow metric would
separately require
\begin{equation}\phantomsection\label{eq-common-lyapunov}{
A_i^\top M+MA_i\preceq-2\lambda M\quad\text{for each mode }i.
}\end{equation} Contracting flows within each mode do not imply this
jump condition. The scalar flow \(\dot x=-x\) with a scheduled reset
\(x^+=2x^-\) every half second has net multiplier \(2e^{-1/2}>1\). Flow
contraction is overwhelmed by repeated expansion at resets. If jumps
have bounded expansion rather than nonexpansion, a dwell-time or
cumulative flow/jump bound may establish an overall result. A common
flow metric alone is insufficient for an impact model.

\section{An Executable Certificate
Example}\label{an-executable-certificate-example}

The following example verifies the double-integrator identity and metric
rate. It is a complete algebraic certificate for the stated linear
model, not a DDP solver or a benchmark of golf performance. The
generalized symmetric eigenvalue calculation avoids explicitly
constructing \(S^{-1/2}\).

\begin{Shaded}
\begin{Highlighting}[]
\ImportTok{import}\NormalTok{ numpy }\ImportTok{as}\NormalTok{ np}
\ImportTok{from}\NormalTok{ scipy.linalg }\ImportTok{import}\NormalTok{ eigvalsh, solve\_continuous\_are}

\NormalTok{A }\OperatorTok{=}\NormalTok{ np.array([[}\FloatTok{0.0}\NormalTok{, }\FloatTok{1.0}\NormalTok{], [}\FloatTok{0.0}\NormalTok{, }\FloatTok{0.0}\NormalTok{]])}
\NormalTok{B }\OperatorTok{=}\NormalTok{ np.array([[}\FloatTok{0.0}\NormalTok{], [}\FloatTok{1.0}\NormalTok{]])}
\NormalTok{Q }\OperatorTok{=}\NormalTok{ np.diag([}\FloatTok{16.0}\NormalTok{, }\FloatTok{1.0}\NormalTok{])}
\NormalTok{R }\OperatorTok{=}\NormalTok{ np.array([[}\FloatTok{1.0}\NormalTok{]])}
\NormalTok{S }\OperatorTok{=}\NormalTok{ solve\_continuous\_are(A, B, Q, R)}
\NormalTok{K }\OperatorTok{=}\NormalTok{ np.linalg.solve(R, B.T }\OperatorTok{@}\NormalTok{ S)}
\NormalTok{A\_closed }\OperatorTok{=}\NormalTok{ A }\OperatorTok{{-}}\NormalTok{ B }\OperatorTok{@}\NormalTok{ K}
\NormalTok{W }\OperatorTok{=}\NormalTok{ Q }\OperatorTok{+}\NormalTok{ K.T }\OperatorTok{@}\NormalTok{ R }\OperatorTok{@}\NormalTok{ K}
\NormalTok{residual }\OperatorTok{=}\NormalTok{ A\_closed.T }\OperatorTok{@}\NormalTok{ S }\OperatorTok{+}\NormalTok{ S }\OperatorTok{@}\NormalTok{ A\_closed }\OperatorTok{+}\NormalTok{ W}
\NormalTok{rate }\OperatorTok{=} \BuiltInTok{float}\NormalTok{(eigvalsh(W, S)[}\DecValTok{0}\NormalTok{] }\OperatorTok{/} \DecValTok{2}\NormalTok{)}
\ControlFlowTok{assert}\NormalTok{ np.linalg.eigvalsh(S)[}\DecValTok{0}\NormalTok{] }\OperatorTok{\textgreater{}} \DecValTok{0}
\ControlFlowTok{assert}\NormalTok{ np.linalg.norm(residual, }\BuiltInTok{ord}\OperatorTok{=}\DecValTok{2}\NormalTok{) }\OperatorTok{\textless{}} \FloatTok{1e{-}10}
\ControlFlowTok{assert}\NormalTok{ np.isclose(rate, }\FloatTok{1.27639320225}\NormalTok{)}
\BuiltInTok{print}\NormalTok{(}\StringTok{"Gain:"}\NormalTok{, np.}\BuiltInTok{round}\NormalTok{(K, }\DecValTok{6}\NormalTok{))}
\BuiltInTok{print}\NormalTok{(}\StringTok{"Metric rate (1/s):"}\NormalTok{, }\BuiltInTok{round}\NormalTok{(rate, }\DecValTok{6}\NormalTok{))}
\end{Highlighting}
\end{Shaded}

The gain is \([4,3]\) and the metric rate is \(1.276393\) inverse
seconds. For a nonlinear model, evaluate the actual policy derivative
and the actual metric derivative. Assigning a Riccati derivative to an
arbitrary guessed matrix at one state does not establish that a positive
metric solves the required differential equation over a trajectory or
region.

A numerical study comparing optimizers must publish model parameters,
coordinate conventions, integrator and step size, input bounds, costs,
initialization, seeds, convergence criteria, metric normalization,
feasibility tolerances, and code revisions. Report failed runs and
distinguish sampled error fits from worst-direction contraction bounds.
Timing requires hardware and solver details. This article supplies no
measured success-rate, basin-volume, robot-tracking, or runtime
comparison; such results require a reproducible experiment package. An
impulse is measured in newton-seconds, whereas a force in newtons
requires its duration and profile to define an impulse.

\section{A Golf Investigation That Connects the
Mechanisms}\label{a-golf-investigation-that-connects-the-mechanisms}

Start with the delivered outcome: clubhead velocity, face orientation,
contact point, and relevant ball-launch variables at the model's own
impact event. A state fixed at a common clock time is a different
comparison from a state evaluated when the club reaches the ball.
Event-time sensitivity becomes important if perturbed swings arrive
earlier or later.

Specify the mechanical state and its history variables. Posture and grip
geometry change load transmission. Inertia couples accelerations among
links. Earlier muscle work and elastic storage affect the available
kinetic energy. Ground reactions can redirect momentum while doing zero
work at an ideal stationary contact. Activation and tissue dynamics set
feasible corrections and impedance. A same-state affine decomposition
describes an acceleration budget; a forward intervention follows a
different state and recomputes those mechanisms.

For a candidate feedback policy, propagate the full variational system
and contact sensitivities. With state-transition matrix \(\Phi(T,t)\)
and a smooth fixed-time outcome \(y_T=H(x(T))\), the first-order output
perturbation is \[
\delta y_T=H_x\Phi(T,t)\delta x(t)
+\int_t^T H_x\Phi(T,s)B(s)\delta u(s)\,ds.
\] For an impact-defined outcome, replace this fixed-time expression by
the appropriate event map and include its timing derivatives. The output
map, remaining time, disturbance direction, feasible control set, and
state uncertainty together determine which errors matter. A metric can
express a useful bound, but a large full-state contraction rate is not
the only route to small impact variability.

Compare candidate mechanisms using matched outcomes and declared
alternatives. A configuration may transmit torque effectively while
amplifying a face-orientation error. Co-contraction may reduce one
perturbation while increasing energetic cost or changing another
response. A low-torque continuation may retain substantial speed without
preserving the required contact state. A golfer may allow harmless joint
variation while correcting a smaller task-relevant component. These
possibilities motivate measurements; none follows solely from an
algebraic decomposition or a fitted exponential.

To compare golfers, use the same outcome definition, model class, units,
uncertainty treatment, feasible input assumptions, and metric
normalization. A computed Riccati matrix depends on chosen cost weights
and policy; it is not a directly measured neural strategy. Do not infer
that skilled golfers maximize contraction, or that a particular local
rate predicts skill, without a preregistered comparison and evidence
against alternative explanations. Useful tests include controlled
perturbations, repeated trials, parameter sensitivity, grip/contact
alternatives, and out-of-sample prediction of delivered impact changes.

\section{Exercises and Checks}\label{exercises-and-checks}

\begin{enumerate}
\def\labelenumi{\arabic{enumi}.}
\tightlist
\item
  Derive the metric derivative and reproduce the matrix-order
  counterexample.
\item
  Compute the alternating discrete product and compare it with its
  individual spectra.
\item
  Solve the finite-horizon scalar example and explain why decreasing
  remaining cost does not imply decreasing state error.
\item
  Verify the double-integrator ARE, spectral rate, and metric rate
  separately.
\item
  Compare the endpoint pullback with the physical two-link mass matrix
  at straight and bent configurations.
\item
  Verify the critically damped error-state certificate and identify why
  its simple energy proof is not strict.
\item
  Compute the scalar noise variance for common and independent noise
  realizations.
\item
  Combine flow and reset multipliers for several event intervals, then
  include event-time sensitivity in a non-grazing guard example.
\item
  Define a golf outcome map and identify one state direction that is
  weakly observed by that outcome. Explain what additional measurements
  would distinguish low output sensitivity from active error correction.
\end{enumerate}

\section*{References and Scope}\label{references-and-scope}
\addcontentsline{toc}{section}{References and Scope}

The contraction definitions follow Lohmiller and Slotine, \emph{On
Contraction Analysis for Non-linear Systems}, Automatica 34(6), 683--696
(1998). The relevant Manchester and Slotine paper is \emph{Control
Contraction Metrics: Convex and Intrinsic Criteria for Nonlinear
Feedback Design}, with the amendment linked above. Riccati and DDP
identities can be checked against the linked MIT chapters; the numerical
counterexamples and certificate code make the distinctions independently
reproducible. The saltation review supplies the event-sensitivity
framework.

These are mathematical and computational sources. They do not validate a
particular golf intervention, neural objective, or physiological
efficiency claim. The contribution here is a consistent chain from model
assumptions through differential equations and feasible controls to
testable outcome predictions.

\textbf{License}: This article and accompanying example code are
released under CC BY 4.0.




\end{document}
